\bookchapter{บทที่ 15: วิทยาศาสตร์แบตเตอรี่กับการยืดอายุการใช้งาน}

\chapterintro{แบตเตอรี่คือ “หัวใจพลังงาน” ของมือถือคุณ บทนี้พาคุณเข้าใจอารมณ์น้อยๆ ของแบต เรียนรู้เทคนิคชาร์จและยืดอายุการใช้งานอย่างมีวิทยาศาสตร์ ลาลังจากความกังวลเรื่องแบตเตอรี่ให้ขาด}

\booksection{หนึ่ง แบตเตอรี่ทำงานยังไง? รอบการชาร์จกับความสุขภาพ}

แบตเตอรี่มือถือ (ลิเธียมไอออนหรือลิเธียมโพลิเมอร์) เหมือนฟองน้ำพลังงาน: ชาร์จคือดูดน้ำ (เก็บไฟ) ใช้งานคือบีบน้ำ (ใช้ไฟ) มีคุณสมบัติน่าจดจำอยู่:

\begin{itemize}
  \item \textbf{ใช้พอชาร์จพอ}: ต่างจากแบตรุ่นเก่า ไม่มีผลความจำ ไม่ต้องรอใช้หมดค่อยชาร์จ
  \item \textbf{กลัวร้อน}: อุณหภูมิสูงคือฆาตกรอันดับหนึ่งของแบต
  \item \textbf{แก่ได้}: ใช้นานๆ ความจุสูงสุดจะลดลง เป็นเรื่องปกติ
\end{itemize}

\booksubsection{ตัวชี้วัดหลัก: จำนวนรอบกับความสุขภาพแบต}

\begin{itemize}
  \item หนึ่งรอบไม่ได้แปลว่าชาร์ดหนึ่งครั้ง เช่นวันนี้ใช้ 60\% ชาร์จเต็ม พรุ่งนี้ใช้ 40\% ชาร์จเต็ม รวมกันครบ 100\% ถึงนับเป็นหนึ่งรอบ
  \item ความสุขภาพแบตคือความจุสูงสุดปัจจุบันเทียบกับตอนใหม่เป็นเปอร์เซ็นต์ ต่ำกว่า 80\% คุณจะรู้สึกได้ชัดว่าอยู่ได้สั้นลง
  \item \textbf{ดูยังไง?}iPhone: ตั้งค่า > แบตเตอรี่ > สุขภาพแบตเตอรี่และการชาร์จ แอนดรอยด์: โดยทั่วไปค้นหา 'สุขภาพแบตเตอรี่' ในการตั้งค่า
\end{itemize}

\booksection{สอง เรื่องลี้ลับการชาร์จ? เร็วไร้สายและกฎทอง}

\begin{itemize}
  \item \textbf{ชาร์จเร็วทำร้ายแบต?}ใช้อุปกรณ์ชาร์จเร็วแท้หรือที่ได้รับการรับรอง ความเสียหายน้อยมาก มือถือกับหัวชาร์จมีระบบควบคุมอุณหภูมิ ร้อนเกินจะลดกำลังชาร์จให้เอง
  \item \textbf{การชาร์จไร้สาย}: สะดวกแต่ประสิทธิภาพต่ำกว่า ร้อนชัดกว่า เหมาะชาร์จทิ้งไว้ทั้งคืน
  \item \textbf{กฎทองดูแลแบต}: พยายามรักษาแบตไว้ระหว่าง 20\%-80\% อย่าให้อด (0\%) หรืออิ่มเกิน (100\%) บ่อยๆ
\end{itemize}

\begin{tcbinfo}{ฟังก์ชัน "จำกัดการชาร์จ 80\%" ของ iPhone — เปลี่ยนกฎทองให้เป็นการตั้งค่าระบบ}
iPhone ซีรีส์ 15/16 (iOS 17 ขึ้นไป) ทำ “กฎ 20\%-80\%” เป็นฟังก์ชันระบบแล้ว เปิดแล้วเครื่องจะหยุดชาร์จเองที่ 80\% ช่วยชะลอความแก่ของแบตได้อย่างมาก เส้นทางตั้งค่า:

\texttt{ตั้งค่า > แบตเตอรี่ > สุขภาพแบตเตอรี่และการชาร์จ > การเพิ่มประสิทธิภาพการชาร์จ > จำกัด 80\%}

ความหมายของสามตัวเลือก: (1)\textbf{ไม่จำกัด}— ชาร์จถึง 100\% (2)\textbf{การชาร์จแบตเตอรี่แบบเพิ่มประสิทธิภาพ}— เรียนรู้จังหวะใช้งานของคุณ ชาร์จให้เต็ม 100\% ช้าลงเฉพาะช่วงท้ายกลางคืน (3)\textbf{จำกัดที่ 80\%}— หยุดที่ 80\% เลย ดูแลแบตที่สุด คนที่อยู่บ้าน/ออฟฟิศมีที่ชาร์จตลอด แนะนำอย่างยิ่งให้เลือกจำกัด 80\% คนที่เดินทางบ่อยต้องแบตเต็มออกเรือน ก็เลือกการชาร์จแบบเพิ่มประสิทธิภาพ

\end{tcbinfo}

\begin{tcbinfo}{หัวชาร์จ GaN — คู่หูสายเดินทาง}
คุณอาจสังเกตว่าหัวชาร์จเร็วตอนนี้เล็กลงเรื่อยๆ ผู้มีพระคุณข้างหลังคือ\textbf{GaN (แกลเลียมไนไตรด์)}วัสดุเซมิคอนดักเตอร์ชนิดใหม่ เทียบกับหัวชาร์จซิลิคอนแบบดั้งเดิม GaN มีข้อดีชัดเจน:

\begin{itemize}
  \item \textbf{ตัวเล็ก}: กำลังเท่ากัน หัวชาร์จ GaN เล็กได้เหลือครึ่งหรือหนึ่งในสามของแบบซิลิคอน หัว GaN 65W เล็กกว่าหัว 5W เก่าของแอปเปิลซะอีก
  \item \textbf{ประสิทธิภาพสูง ร้อนต่ำ}: การสูญเสียที่สวิตชิงของแกลเลียมไนไตรด์ต่ำกว่าซิลิคอนมาก ประสิทธิภาพการแปลงพลังงานเกิน 95\% ตอนชาร์จไม่ค่อยร้อนมือ
  \item \textbf{เหมาะไปเดินทาง}: หัวชาร์จ GaN หลายช่อง 65W หรือ 100W ชาร์จมือถือ แท็บเล็ต โน้ตบุ๊กพร้อมกันได้ สบายกว่าพกก้อน "อิฐ" สามก้อนเยอะ
  \item \textbf{คำแนะนำการเลือกซื้อ}: เลือกหัวชาร์จ GaN ที่รองรับโปรโตคอล PD (65W ขึ้นไปชาร์จโน้ตบุ๊กได้) แบรนด์ที่รู้จักอย่าง Baseus, Anker, Nubia, Xiaomi ฯลฯ ใช้ได้หมด อย่าลืมเช็กว่าสายข้อมูลเข้ากันไหม (กำลังสูงต้องใช้สาย Type-C ที่มีชิป E-Marker)
\end{itemize}

\end{tcbinfo}

คุณอาจเคยเจอสถานการณ์นี้: มือถือรองรับชาร์จเร็วชัดๆ แต่เปลี่ยนหัวชาร์จไปกลับช้าลง เพราะชาร์จเร็วไม่ใช่แค่ “กำลังใหญ่พอ” แต่ต้องให้หัวชาร์จ สายข้อมูล มือถือ สามฝ่าย “พูดภาษาเดียวกัน” — นั่นคือโปรโตคอลชาร์จเร็ว

\booksubsection{หลักการพื้นฐานของชาร์จเร็ว}

กำลังชาร์จ = แรงดัน × กระแส (P = U × I) จะเพิ่มกำลังก็ต้องเพิ่มแรงดัน เพิ่มกระแส หรือเพิ่มทั้งคู่ แต่ละโปรโตคอลเลือกเส้นทางเทคโนโลยีต่างกัน

ขั้นตอนการทำงานของชาร์จเร็ว:

\begin{enumerate}
  \item \textbf{ขั้นจับมือ (Handshake)}: มือถือกับหัวชาร์จสื่อสารกันผ่านสายข้อมูล แลกเปลี่ยนโปรโตคอลและข้อมูลกำลังที่แต่ละฝ่ายรองรับ
  \item \textbf{ขั้นเจรจา}: ทั้งสองฝ่ายกำหนดระดับกำลังสูงสุดที่รองรับร่วมกัน
  \item \textbf{ขั้นชาร์จ}: หัวชาร์จส่งกระแสแรงดันตามที่เจรจาไว้ มือถือทำการจัดการพลังงานฝั่งตัวเอง
  \item \textbf{ขั้นปรับแต่ง}: พอแบตเพิ่มขึ้น มือถือจะลดกำลังชาร์จทีละขั้นเพื่อป้องกันแบตเตอรี่
\end{enumerate}

\booksubsection{โปรโตคอลชาร์จเร็วกระแสหลักฉบับละเอียด}

\textbf{USB PD (Power Delivery)}— ราชาแห่งโปรโตคอลสากล

\begin{itemize}
  \item \textbf{ลักษณะเด่น}: มาตรฐานทางการที่ USB-IF กำหนด ครอบคลุมกว้างที่สุด PD 3.0 รองรับสูงสุด 100W (20V/5A) PD 3.1 ขยายถึง 240W (48V/5A)
  \item \textbf{ระดับแรงดัน}: 5V, 9V, 15V, 20V PD 3.1 ยังรองรับ 28V, 36V, 48V
  \item \textbf{ความเข้ากันได้}: iPhone ทุกรุ่น Google Pixel มือถือฟลากชิปแอนดรอยด์บางรุ่น โน้ตบุ๊ก เครื่องเกม Switch ล้วนรองรับ PD
  \item \textbf{ข้อดี}: ใช้สากล หัวชาร์จเดียวชาร์จได้หลายอุปกรณ์ รองรับ PPS (แหล่งจ่ายไฟโปรแกรมได้) ปรับแรงดันกระแสได้ละเอียด
  \item \textbf{ข้อเสีย}: ต้องใช้หัวเสียบ Type-C และชิปเฉพาะ ต้นทุนสูงกว่า
\end{itemize}

\textbf{Qualcomm QC (Quick Charge)}— อดีตกาลฝ่ายแอนดรอยด์

\begin{itemize}
  \item \textbf{QC 3.0}: รองรับแรงดันปรับได้ 3.6V-20V ช่วงละ 0.2V สูงสุด 18W
  \item \textbf{QC 4.0/4+}: เข้ากันได้กับโปรโตคอล USB PD สูงสุด 27W เพิ่มเทคโนโลยีชาร์จสองวงจร
  \item \textbf{QC 5.0}: รองรับสูงสุด 100W+ 0-50\% ใช้เวลาแค่ 5 นาที ปรับปรุงความร้อนได้มาก
  \item \textbf{ความเข้ากันได้}: มือถือที่ใช้ชิป Snapdragon ส่วนใหญ่รองรับ แต่แต่ละแบรนด์มักใช้โปรโตคอลส่วนตัวของตัวเองเป็นหลัก
\end{itemize}

\textbf{OPPO VOOC/OnePlus WARP/realme Dart}— ตัวแทนแรงดันต่ำกระแสสูง

\begin{itemize}
  \item \textbf{เส้นทางเทคโนโลยี}: แรงดันต่ำกระแนสูง (เริ่มที่ 5V/4A) ร้อนน้อยกว่า ประสิทธิภาพสูงกว่า
  \item \textbf{VOOC 1.0}: 20W (5V/4A) เปิดตัวปี 2014 เปิดยุคชาร์จเร็ว
  \item \textbf{SuperVOOC}: 50W-65W ใช้ดีไซน์เซลล์คู่ต่ออนุกรม
  \item \textbf{SuperVOOC 2.0}: 65W ชาร์จแบต 4000mAh เต็มใน 30 นาที
  \item \textbf{ซูเปอร์แฟลชชาร์จ 150W/240W}: เทคโนโลยีชาร์จเร็วที่ระดับการผลิตเร็วที่สุดตอนนี้
  \item \textbf{ลักษณะเด่น}: ต้องใช้หัวชาร์จและสายข้อมูลเฉพาะ (มีขั้วต่อเพิ่ม) เข้ากันไม่ได้กับโปรโตคอลอื่น
\end{itemize}

\textbf{Huawei SuperCharge}— วิวัฒนาการชาร์จเร็วแรงดันสูง

\begin{itemize}
  \item \textbf{FCP (Fast Charge Protocol)}: 9V/2A, 18W เข้ากันได้กับ QC 2.0
  \item \textbf{SCP (Super Charge Protocol)}: รองรับกำลังสูงขึ้น จาก 22.5W สู่ 66W และ 100W
  \item \textbf{จุดเด่นเทคโนโลยี}: ใช้เทคโนโลยีปั๊มประจุ แปลงแรงดันสูงเป็นกระแนสูงแรงดันต่ำเข้าสู่แบต ประสิทธิภาพถึง 98\%
  \item \textbf{ความเข้ากันได้}: เฉพาะ Huawei/Honor หัวชาร์จบางรุ่นรองรับ PD
\end{itemize}

\textbf{Xiaomi Surge ชาร์จเร็ว}— การรุกของโปรโตคอลพัฒนาเอง

\begin{itemize}
  \item \textbf{ชิป Surge P1}: ชิปจัดการชาร์จเร็วพัฒนาเองของ Xiaomi รองรับชาร์จเร็วมีสาย 120W
  \item \textbf{ชิป Surge G1}: ชิปจัดการแบตเตอรี่ ทำงานคู่ P1 ให้ชาร์จเร็วมีประสิทธิภาพยิ่งขึ้น
  \item \textbf{จุดเด่นเทคโนโลยี}: ใช้เซลล์คู่ต่ออนุกรม + ปั๊มประจุ รองรับโปรโตคอลหลายตัวได้
\end{itemize}

\textbf{ชาร์จเร็วแอปเปิล}— ระมัดระวังแต่มั่นคง

\begin{itemize}
  \item \textbf{โปรโตคอล}: รองรับเฉพาะโปรโตคอล USB PD
  \item \textbf{กำลัง}: iPhone 8 ขึ้นไปรองรับชาร์จเร็ว สูงสุดราว 27W (iPhone 15 Pro Max ถึง 27W+)
  \item \textbf{ลักษณะเด่น}: กำลังค่อนข้างระมัดระวัง แต่การจัดการเส้นโค้งการชาร์จดี ควบคุมความร้อนได้เยี่ยม
  \item \textbf{ข้อกำหนดอุปกรณ์เสริม}: ต้องใช้สาย USB-C to Lightning (หรือ USB-C to USB-C) และหัวชาร์จ PD
\end{itemize}

\booksubsection{ตารางเทียบโปรโตคอลชาร์จเร็ว}

\begin{center}
\begin{tabularx}{\linewidth}{|X|X|X|X|}
\hline
\textbf{โปรโตคอล} & \textbf{กำลังสูงสุด} & \textbf{เส้นทางเทคโนโลยี} & \textbf{ความเข้ากันได้} \tabularnewline
\hline
USB PD 3.1 & 240W & แรงดันสูง/หลายระดับ & กว้างขวางที่สุด \tabularnewline
QC 5.0 & 100W+ & แรงดันสูงปรับได้ & เครื่อง Qualcomm \tabularnewline
VOOC/Dart & 240W & แรงดันต่ำกระแนสูง & เฉพาะสาย OPPO \tabularnewline
SuperCharge & 100W & ปั๊มประจุ & Huawei/Honor \tabularnewline
ชาร์จเร็วแอปเปิล & \textasciitilde{}27W & โปรโตคอล PD & iPhone/iPad \tabularnewline
\hline
\end{tabularx}
\end{center}

\booksubsection{โปรโตคอลชาร์จไร้สาย}

การชาร์จไร้สายก็มีมาตรฐานโปรโตคอลของตัวเอง:

\begin{itemize}
  \item \textbf{มาตรฐาน Qi}: กำหนดโดย WPC (สมาคมชาร์จไร้สาย) มาตรฐานที่ใช้สากลที่สุด Qi พื้นฐานรองรับ 5W-15W
  \item \textbf{มาตรฐาน Qi2}: เปิดตัวปี 2023 เพิ่มการจัดตำแหน่งแม่เหล็ก (คล้าย MagSafe ของแอปเปิล) รองรับ 15W+
  \item \textbf{โปรโตคอลส่วนตัว}: แต่ละแบรนด์มีชาร์จไร้สายความเร็วสูงของตัวเอง เช่น Xiaomi 50W Huawei 50W ฯลฯ ต้องใช้ที่ชาร์จเฉพาะ
\end{itemize}

\booksubsection{คู่มือเลือกซื้อชาร์จเร็ว}

\begin{itemize}
  \item \textbf{ดูความตรงกันของโปรโตคอล}: ซื้อหัวชาร์จให้ยืนยันว่ารองรับโปรโตคอลชาร์จเร็วของมือถือคุณ หัวชาร์จภายนอกมักติดป้ายบอกโปรโตคอลที่รองรับ
  \item \textbf{ดูสายข้อมูล}: ชาร์จเร็วกำลังสูงต้องใช้สายเฉพาะ สายทั่วไปอาจรองรับแค่ต่ำกว่า 60W
  \item \textbf{ดูการรับรอง}: ผู้ใช้แอปเปิลแนะนำเลือกสายรับรอง MFi ฝั่งแอนดรอยด์เลือกของแท้แบรนด์หรือภายนอกที่มีชื่อเสียง
  \item \textbf{ดูความร้อน}: ตอนชาร์จเร็วมือถือร้อนขึ้นปกติ แต่ถ้าร้อนถึงมือจนไม่ถนัดไม่ปกติ แนะนำหยุดชาร์จหรือเปลี่ยนหัวชาร์จ
\end{itemize}

\begin{tcbwarn}{คำเตือนความปลอดภัยชาร์จเร็ว}
\begin{itemize}
  \item เลี่ยงเล่นเกมใหญ่ไปพลางชาร์จไปพลาง จะยิ่งร้อนกว่าเดิม
  \item อย่าใช้หัวชาร์จและสายข้อมูลคุณภาพต่ำ อาจก่อความเสี่ยงด้านความปลอดภัย
  \item ตอนชาร์จถอดเคสหนาออก ช่วยระบายความร้อน
  \item ถ้าหัวชาร์จหรือมือถือร้อนผิดปกติ มีกลิ่นแปลก หยุดใช้ทันที
\end{itemize}

\end{tcbwarn}

\booksection{สี่ เทคโนโลยีแบตใหม่: แบตแอโนดซิลิคอนคาร์บอน}

ถ้าคุณตามงานเปิดตัวมือถือฟลากชิปปี 2024-2026 คำว่า “แอโนดซิลิคอนคาร์บอน” แทบกลายเป็นคำราชาของแบตระดับสูง นี่คือการยกระดับสำคัญของวัสดุแอโนดลิเธียม หลักการจริงๆ ไม่ซับซ้อน:

\begin{itemize}
  \item \textbf{แอโนดกราไฟต์แบบดั้งเดิม}: แบตมือถือใช้กราไฟต์เป็นแอโนดมาตลอด เทคโนโลยีโตแล้ว มั่นคง แต่ความหนาแน่นพลังงานใกล้ขีดจำกัดทางทฤษฎีแล้ว (ราว 372mAh/g)
  \item \textbf{ผสมวัสดุซิลิคอน}: ความจุจำเพาะทางทฤษฎีของซิลิคอนสูงกว่ากราไฟต์ราว 10 เท่า (ราว 4200mAh/g) ผสมซิลิคอนน้อยๆ ลงในกราไฟต์ทำเป็น "แอโนดคอมโพสิตซิลิคอนคาร์บอน" ช่วยเพิ่มความหนาแน่นพลังงานของแบตทั้งก้อนได้ชัดเจน
  \item \textbf{ความหนาแน่นพลังงานเพิ่มขึ้น}: ขนาดเท่ากัน แบตแอโนดซิลิคอนคาร์บอนความจุเพิ่มได้ราว\textbf{15\%-20\%}นี่คือเหตุผลสำคัญที่มือถือฟลากชิปปี 2024-2026 ยัดแบตใหญ่ 5500-6000mAh ลงตัวเครื่องบาง 7-9mm ได้
\end{itemize}

\begin{tcbinfo}{รุ่นไหนใช้แล้วบ้าง?}
\begin{itemize}
  \item \textbf{Honor}: ซีรีส์ Magic เป็นรายแรกที่ใช้เชิงพาณิชย์ Magic6 Pro, ซีรีส์ Magic7, Magic V3 พับได้ ล้วนใช้
  \item \textbf{Xiaomi}: Xiaomi 14 Ultra, ซีรีส์ Xiaomi 15 ติดตั้งแบตแอโนดซิลิคอนคาร์บอนรุ่นใหม่
  \item \textbf{vivo}: X100 Ultra, ซีรีส์ X200 Pro ใช้แผนแอโนดซิลิคอนคาร์บอนความหนาแน่นสูง
  \item \textbf{Huawei/OPPO ฯลฯ}: ปี 2025-2026 ฟลากชิปตามรอยกันทั่วไป แอโนดซิลิคอนคาร์บอนกำลังเปลี่ยนจาก "เทคดำ" เป็นอุปกรณ์มาตรฐานของเครื่องระดับสูง
\end{itemize}

\end{tcbinfo}

\textbf{ผลกระทบจริงต่อผู้ใช้}: ประโยชน์ใหญ่สุดของแอโนดซิลิคอนคาร์บอนคือ “ขนาดเท่ากันแต่ความจุแบตมากขึ้น” — เครื่องทำบางลงได้พร้อมอยู่ได้ไม่สั้นลง หรือไม่หนาขึ้นเลยแต่แบตเด้งจาก 4500mAh สู่ 5500mAh กล่าวอีกนัย คุณไม่ต้องจ่ายค่า “หนาหนัก” เพื่อความอึดแล้ว ต้องระวังว่าแบตแอโนดซิลิคอนคาร์บอนยังต้องปรับจูนกับชิปชาร์จเร็วในเรื่องอุณหภูมิต่ำและการจัดการชาร์จเต็ม นี่ก็เป็นเหตุผลหนึ่งที่ผู้ผลิตแนะนำให้คุมการชาร์จไว้ที่ 80\% ระยะยาว

\booksection{ห้า นี่แหละเคล็ดลับประหยัดแบตตัวจริง}

\begin{itemize}
  \item \textbf{หน้าจอ (ราชากินแบต)}: ลดความสว่าง ลดเวลาล็อกจออัตโนมัติ จอ OLED เปิดโหมดมืดช่วยประหยัดไฟได้เยอะ
  \item \textbf{สัญญาณอ่อน}: ในลิฟต์ ใต้ดิน มือถือจะพยายามหาสัญญาณสุดชีวิต กินไฟรวดเร็วมาก
  \item \textbf{5G}: ถ้าสัญญาณไม่ดี สลับกลับ 4G ด้วยมือ อยู่ได้นานขึ้นเยอะ
  \item \textbf{บริการระบุตำแหน่ง}: แอปที่ไม่ใช่พวกนำทาง ตั้งสิทธิ์ตำแหน่งเป็น "อนุญาตเฉพาะขณะใช้งาน"
  \item \textbf{จำกัดกิจกรรมพื้นหลัง}: iPhone ที่ ตั้งค่า > ทั่วไป > การรีเฟรชแอปในพื้นหลัง แอนดรอยด์ที่การตั้งค่าแบตของแอป
\end{itemize}

\begin{tcbpractice}{มาตรวจดู: ทำการตรวจสุขภาพให้แบตของคุณ}
\begin{enumerate}
  \item ดูความสุขภาพแบต: หาและจดความจุสูงสุดของแบตคุณไว้
  \item วิเคราะห์รายละเอียดการกินไฟ: ดูกราฟกับรายการแอปในการตั้งค่าแบต หาตัวร้ายตัวกินไฟ
  \item วางแผนรับมือ: จำกัดแอปที่กินไฟผิดปกติในพื้นหลัง
\end{enumerate}

\end{tcbpractice}
